\documentclass[../../main.tex]{subfiles}

\begin{document}

\section{$C(X)$ 的对偶与 Riesz-Kakutani 表示定理}\label{section:ch21-rieszkakutani}

令 $X$ 为紧 Hausdorff 空间. 我们赋予 $C(X)$ 以最大值范数
\begin{align*}
\|f\|_{\max}=\max_{x\in X}|f(x)|.
\end{align*}
我们希望刻画 Banach 空间 $C(X)$ 的对偶. 根据 Riesz-Markov 表示定理, $C(X)$ 上的正线性泛函由 Radon 测度的积分给出. 然而, $C(X)$ 上的一般连续线性泛函不是正的. 我们首先证明每个连续线性泛函可表示为两个正线性泛函的差.

\begin{definition}[正部分与负部分]\label{definition:ch21-正负部分}
令 $\Lambda$ 为 $C(X)$ 上的连续线性泛函. 对 $f\geqslant 0$ 定义
\begin{align*}
\Lambda^{+}(f)=\sup\{\Lambda(g)\mid g\in C(X),\ 0\leqslant g\leqslant f\}.
\end{align*}
然后通过线性延拓到所有 $C(X)$. 定义 $\Lambda^{-}=\Lambda^{+}-\Lambda$.
\end{definition}

\begin{lemma}\label{lemma:ch21-正负部分正泛函}
$\Lambda^{+}$ 和 $\Lambda^{-}$ 都是 $C(X)$ 上的正线性泛函, 且 $\Lambda=\Lambda^{+}-\Lambda^{-}$.
\end{lemma}

\begin{proof}
显然 $\Lambda^{+}(f)\geqslant 0$ 对 $f\geqslant 0$. 线性性的验证是标准的. $\Lambda^{-}$ 的正性来自 $\Lambda^{+}(f)\geqslant\Lambda(f)$ 对 $f\geqslant 0$.
\end{proof}

\begin{theorem}[Riesz-Kakutani 表示定理]\label{theorem:ch21-rieszkakutani}
令 $X$ 为紧 Hausdorff 空间而 $\Lambda$ 是 $C(X)$ 上的有界线性泛函. 则存在 $\mathcal{B}(X)$ 上的唯一带号 Radon 测度 $\nu$ 使得
\begin{align}\label{align:ch21-rieszkakutani等式}
\text{对所有 }f\in C(X),\quad\Lambda(f)=\int_{X}f\,d\nu.
\end{align}
此外, $\|\Lambda\|=|\nu|(X)$, 其中 $|\nu|$ 是 $\nu$ 的全变差测度.
\end{theorem}

\begin{proof}
根据前一个引理, $\Lambda=\Lambda^{+}-\Lambda^{-}$, 其中 $\Lambda^{+}$ 和 $\Lambda^{-}$ 是正线性泛函. 根据 Riesz-Markov 定理, 存在 Radon 测度 $\mu^{+}$ 和 $\mu^{-}$ 使得
\begin{align*}
\text{对所有 }f\in C(X),\quad\Lambda^{+}(f)=\int f\,d\mu^{+},\quad\Lambda^{-}(f)=\int f\,d\mu^{-}.
\end{align*}
定义带号测度 $\nu=\mu^{+}-\mu^{-}$. 则\eqref{align:ch21-rieszkakutani等式}成立. 唯一性来自 Radon-Nikodym 定理或直接的测度论论证.

范数等式 $\|\Lambda\|=|\nu|(X)$ 的证明: 一方面,
\begin{align*}
|\Lambda(f)|=\left|\int f\,d\nu\right|\leqslant\|f\|_{\max}\cdot|\nu|(X),
\end{align*}
因此 $\|\Lambda\|\leqslant|\nu|(X)$. 另一方面, 根据带号测度的 Hahn 分解, 存在集合 $P$ 使得 $\nu^{+}(E)=\nu(E\cap P)$ 且 $\nu^{-}(E)=-\nu(E\cap P^{C})$. 令 $f$ 为逼近 $\chi_{P}-\chi_{P^{C}}$ 的连续函数. 则 $\Lambda(f)\approx\nu(P)-\nu(P^{C})=|\nu|(X)$. 因此 $\|\Lambda\|\geqslant|\nu|(X)$.
\end{proof}

我们称 $C(X)$ 的对偶中的元素为 $X$ 上的有号 Radon 测度. 该对偶空间记为 $M(X)$ (测度空间).

\begin{corollary}\label{corollary:ch21-弱星紧}
令 $X$ 为紧 Hausdorff 空间. 则 $M(X)$ 中的闭单位球是弱 $*$ 紧的.
\end{corollary}

\begin{proof}
这是 Alaoglu 定理的直接推论.
\end{proof}

\subsection*{习题}

\begin{exercise}\label{exercise:ch21-31}
验证 $\Lambda^{+}$ 的线性性.
\end{exercise}

\begin{exercise}\label{exercise:ch21-32}
证明 Riesz-Kakutani 定理中的唯一性.
\end{exercise}

\begin{exercise}\label{exercise:ch21-33}
令 $\nu$ 为带号 Radon 测度. 证明它的全变差 $|\nu|$ 也是 Radon 测度.
\end{exercise}

\begin{exercise}\label{exercise:ch21-34}
用 Krein-Milman 定理证明: $M(X)$ 的单位球的极点对应于 $X$ 上的概率测度.
\end{exercise}

\end{document}